
@@M0@@
Pour éviter toute confusion, désignons par @@M1@@ la classe de @@M2@@ dans @@M3@@. Si @@M4@@ est l'immersion canonique, alors la formule de projection (IV 2.11.1.2) donne
@@M5@@
On a @@M6@@, donc d'après (1.1.8) et (1.3.3) pour @@M7@@, @@M8@@, on trouve
@@M9@@
d'où l'assertion.

Le théorème 1.3.2 est trivial pour @@M10@@. Supposons-le démontré par récurrence pour @@M11@@ et pour tout schéma noethérien @@M12@@. D'après la remarque précédente il nous reste à démontrer que pour tout schéma noethérien \emph{intègre} @@M13@@ de dimension @@M14@@ et pour tout produit (1.2.3) satisfaisant à (1.2.4), on a
@@M15@@
Si @@M16@@ dans (1.2.3), on a gagné grâce à l'hypothèse de récurrence et l'associativité de la loi (1.3.1). On peut donc supposer (@@M17@@ étant connexe) que
@@M18@@
où @@M19@@ et @@M20@@ est un @@M21@@-module localement libre de rang @@M22@@.

Distinguons deux cas :

I. Supposons qu'on peut trouver des @@M23@@-modules inversibles @@M24@@ tels que
@@M25@@
